\section{Conclusion}
\label{sec:conclusion}
We presented a capability-matched test for reasoning-LLM distillation provenance. A
KL decomposition shows that the natural base-relative likelihood surplus conflates a
candidate's capability with distillation descent, explaining why it false-alarms on
honestly trained students and is blind to same-family distillation. Replacing the
weak base model with a same-base non-distilled reference scored under the same
candidate, $r_c(S)=\overline{\ell_c(S)-\ell_c(R)}$, controls the capability term to
first order and isolates descent. The matched trace-level test detects distillation
at AUROC $\geq0.99$ in the three-seed GSM8K grid and in a single-seed MATH stress
test, recovering same-family distillation that the natural statistic misses
entirely (AUROC $0.13\!\to\!0.998$ on GSM8K). With MTCR, closed-set R1-Qwen-32B and
R1-Llama-8B attribution rises from $0.658$--$0.672$ to $0.875$--$0.894$, while
shared-ancestor detection remains high.

The revised stress summaries push the paper beyond the original clean-trace audit:
CML+MTCR retains AUROC at least $0.930$ under the reported adaptive
trace-transformation sweeps, the cross-base/task matrix stays above AUROC $0.912$,
and open-set false attribution remains at or below $2.0\%$ in the reported decoy
conditions. These results support a sharper TIFS-style claim: a non-invasive,
post-hoc forensic audit can control capability false positives, recover same-family
distillation, and remain useful under several realistic reviewer stress tests. The
main operating assumptions are clear: MTCR uses the candidate set and calibration
views defined by the audit, and CML uses a known or declared same-base reference.
The broader lesson is simple: provenance should be
measured against a capability-matched reference, not against a weak absolute
baseline.
